\documentclass[10pt,a4paper]{amsart}
\usepackage[margin=19mm]{geometry}
\usepackage{amsmath,amssymb,mathtools}
\usepackage[scr=rsfs,cal=euler]{mathalfa}
\usepackage{xfrac}
\usepackage[unicode,hidelinks,bookmarks=false]{hyperref}
\newcommand{\ie}{i.e.\ }
\newcommand{\cf}{cf.\ }
\newcommand{\resp}{respectively\ }
\newcommand{\Subsection}{\subsection}
\providecommand{\nobreakdash}{\nobreak\hbox{-}}
\setlength{\emergencystretch}{2em}
\multlinegap0pt
\usepackage{bm}
\usepackage{bbm}
\usepackage{dsfont}
\usepackage{xspace}
\usepackage{cancel}
\usepackage{mathtools}
\usepackage{amsthm}
\makeatletter
\@ifundefined{equation}{\newtheorem{equation}{Equation}}{}
\@ifundefined{proposition}{\newtheorem{proposition}[equation]{Proposition}}{}
\makeatother
\def\Z{{\mathbb{Z}}}
\def\epsilon{\varepsilon}
\def\zp{{\Z_p}}
\newcommand{\mbb}[1]{\mathbb{#1}}
\title[p-adic interpolation of Gauss--Manin connections on nearly o]{p-adic interpolation of Gauss--Manin connections on nearly overconvergent modular forms and p-adic L-functions}
\author{Andrew Graham, Vincent Pilloni,, Joaqu\'in Rodrigues Jacinto}
\date{Source: arXiv:2311.14438v2 (CC BY 4.0)}
\begin{document}
\maketitle

\noindent\emph{Source and licence.} p-adic interpolation of Gauss--Manin connections on nearly overconvergent modular forms and p-adic L-functions. Version arXiv:2311.14438v2 (CC BY 4.0), distributed under the Creative Commons Attribution 4.0 International licence, as recorded in the arXiv licence metadata for this exact version and shown on the abstract page https://arxiv.org/abs/2311.14438v2. Full rights notice: licenses/arxiv-2311.14438-CC-BY-4.0.txt.\par\smallskip

\noindent\textit{Editorial scope.} Counted unit: the Proposition whose source label is \texttt{Lemmaepsan1} in the article (source0.tex lines 951--953, source label \texttt{Lemmaepsan1}), delivered together with the article's own complete original proof of it (source0.tex lines 955--984, one \texttt{proof} environment of 30 lines). No mathematics is written, repaired or re-derived: statement and proof are the article's own text line for line. Required context retained uncounted: none. Every definition, display and label of those lines is retained unchanged. Omitted from this package and not redistributed: 1--950, 954--954, 985--2640. Nothing inside a retained excerpt is omitted or altered: every retained body is delivered exactly as the article's lines are written, so no omission receipt of any kind accompanies this package. Citations: the retained excerpts carry no citation command at all, so no bibliography entry of the article is transcribed and no reference is redistributed. Reference adaptation: none was needed and none was made; every reference inside the delivered unit resolves inside it, so no unresolved reference or citation is produced. Printed slips kept literal: no printed word, symbol, hypothesis or number of the counted statement or of its proof was altered. Layout and class: the article's own class is \texttt{amsart} with packages including geometry, amsmath, amssymb, tikz, tikz-cd, todonotes, mathrsfs, extarrows, graphicx, hyperref, IEEEtrantools, inputenc, babel, comment; the wrapper is the lane's pinned \texttt{amsart} editorial wrapper at 10pt on A4, which restates the theorem-like environments the retained text needs and adds the article's own macro definitions that the retained text needs (\texttt{\textbackslash Z}, \texttt{\textbackslash epsilon}, \texttt{\textbackslash mbb}, \texttt{\textbackslash zp}), with extra wrapper packages (bm, bbm, dsfont, xspace, cancel, mathtools). The delivered numbering is the unit's own. Assets: no retained span contains a figure environment, a graphics-inclusion command, a PDF-page inclusion, a table or a TikZ picture; no image file is copied into this package, the package declares no asset dependency and no image file is redistributed with it. Licence: version arXiv:2311.14438v2 of this article is distributed under the Creative Commons Attribution 4.0 International licence, as recorded in the arXiv licence metadata for this exact version and shown on the abstract page https://arxiv.org/abs/2311.14438v2. A full rights notice is delivered at licenses/arxiv-2311.14438-CC-BY-4.0.txt. Characters: every retained excerpt is pure ASCII, so the delivered engine, XeLaTeX with its default Latin Modern text font, reports no missing character.
\begin{proposition} \label{Lemmaepsan1}
Let $T_1, T_2 \in \mathrm{End}_R(V)$ and suppose that $T_1$ extends to a locally analytic action and that $|\!|T_2|\!| \leq 1$. Then for any $\epsilon > 0$, there exists an integer $N_\epsilon \geq 1$ (depending on $T_1$ but not on $T_2$) such that $T_1 + p^{N_\epsilon} T_2$ extends to an $\epsilon$-analytic action.  
\end{proposition}

\begin{proof}
Let $N \geq 0$ and denote $T = T_1 + p^N T_2$. Let $f_k(X)$ denote the polynomial
\[ f_k(X) = \frac{X (X-1) \ldots (X - k + 1)}{k!}. \]
Let $\Lambda$ denote the set of all ordered tuples $I = (-1 = t_0 \leq t_1 < t_2 < \ldots < t_{2r-1} \leq t_{2r} = k-1)$ with $r \geq 1$. For such an $I \in \Lambda$ and any $1 \leq i \leq r$, let $k_i = t_{2i-1} - t_{2i-2}$ and $\ell_i = t_{2i} - t_{2i - 1}$, and set
 \[ z_I = f_{k_1}(T_1 - (t_0 + 1)) \cdot T_2^{\ell_1} \cdot f_{k_2}(T_1 - (t_2 + 1)) \cdot T_2^{\ell_2} \cdot \ldots \cdot f_{k_r}(T_1 - (t_{2r} + 1)) \cdot T_2^{\ell_r} .\]
Then we can write
\begin{equation} \label{eqtmp2alt}
f_k(T) = f_k(T_1) + \sum_{I \in \Lambda, b \neq 0} \frac{p^{N b}}{b!} \binom{k}{a}^{-1} \binom{a}{k_1, \ldots, k_r}^{-1} z_I,
\end{equation}
where we have denoted $a = \sum_{i=1}^r k_i$ and $b = \sum_{i=1}^r \ell_i$ so that $a + b = k$.

We now give a bound for the operators $z_I$. Note that for any integer $t \in \Z$, $f_k(X - t)$ is also a polynomial in $X$ and the Mahler expansion of $f_k(X - t)$ is of the form
\[ f_k(X - t) = \sum_{i = 0}^k a_i f_i(X) \] 
for $a_i \in \Z_p$ because $f_k(x - t) \in \zp$ for any $x \in \zp$ and the coefficients are just computed
using the discrete difference operator. This implies that, if $\epsilon > 0$ and $C \in \mbb{R}_{> 0}$ is such that
\[ p^{-k \epsilon / 2} |\!| f_k(T_1) |\!| \leq C \]
for all $k \geq 0$ then,
for any $t \in \Z$ and $k \geq 0$, we have
\[
p^{-k \epsilon / 2} |\!|f_k(T_1 - t)|\!| \leq C.
\]
Hence, by the fact that $|\!|T_2|\!| \leq 1$, we deduce
\begin{equation} \label{eqtmp1alt} p^{- a \epsilon / 2} |\!|z_I|\!| \leq C^{r}.
\end{equation}
Using that for any multinomial coefficient we have $v_p( \binom{n}{a_1, \ldots, a_s}^{-1} ) \geq - \log_p(n)$, we get from \eqref{eqtmp2alt} and \eqref{eqtmp1alt} that
\[ p^{-k \epsilon} |\!|f_k(T) |\!| \leq \mathrm{max}(C, p^{-k \epsilon} p^{-Nb + b/(p-1)} p^{2 \log_p(k)} p^{(k - b) \epsilon / 2} C^r ). \]
Since $b \geq r - 1$, taking $N$ such that $p^{-N + 1/(p-1)} \leq C^{-1}$ we deduce that $p^{-Nb + b/(p-1)} C^r \leq C$ and the right hand side is bounded above by
\[ p^{-k \epsilon / 2 + 2 \log_p(k)} C \]
which is uniformly bounded in $k$ as $\epsilon > 0$. This concludes the proof.
\end{proof}

\end{document}
